\section{Theoretical Analysis}
\label{section:analysis}
In this section, we establish the following rigorous theoretical results for the proposed numerical method: 
(a) well-posedness and discrete energy stability to the convex splitting scheme;
(b) unconditional linear convergence and bound-preserving property for the ADMM solver.
\subsection{Well-posedness and discrete energy stability}
We first prove the existence and uniqueness of the discrete solution of the convex splitting scheme 
\eqref{scheme:cs} at each time step.
\begin{theorem}\label{thm:well-posedness}
    Assume that $0<u^n<1$ and $u^n \in \ch$. For any time step size $\tau>0$, there exists a unique $u^{n+1}\in \ch$ to solve \eqref{scheme:cs}.
\end{theorem}
\begin{proof}
    From \eqref{eq:opt}, the proof reduces to establishing the existence and uniqueness of a minimizer. 
    The Hessian matrix of the objective function is given by
    \begin{equation}
        \nabla^2 J(u)=\frac{1}{\tau}I-\epsilon^2\Delta_h+G
    \end{equation}
    where $G=\text{diag}\{
    \frac{1}{u_j^n(1-u_j^n)}\}_{j=0}^{N^2-1}$. 
    Noticing the objective function $J(u)$ is $\frac{1}{\tau}$-strongly convex because of the property that $\nabla^2J(u)-\frac{1}{\tau}I$ is always positive-definite, 
    it suffices to show that a minimizer exists; uniqueness is then a direct consequence of strong convexity \cite{boyd2004convex}.
    \par Existence. If one component $u_i\to 0^+$ or $u_i \to 1^-$, the logarithmic term $\log(u_i)-\log(1-u_i)$ tends to $-\infty$ or $+\infty$ while other terms are all finite. 
    Therefore, the $i$-th equation in \eqref{scheme:cs} cannot hold under the assumption. 
    We thus conclude the existence of a positive constant $\delta>0$ such that $u_i^n\in[\delta,1-\delta]$ for all $0\leq i\leq N^2-1$. 
    A direct application of the extreme value theorem to the continuous function $J(u)$ on the compact constraint set $[\delta,1-\delta]^{N^2}$ 
    readily yields the existence of at least one minimizer \cite{1976Principles,weierstrass1894mathematische}.
\end{proof}


For $u\in \ch$, the discrete energy is defined by
\begin{equation}
\label{eq:discrete-energy}
    E_h(u)=\langle F(u),1 \rangle_h+\frac{\epsilon^2}{2}\|\nabla_h u\|_h^2,
\end{equation}
where $F(u)=u\log(u)+(1-u)\log(1-u)+\theta(u-u^2)$. Analogous to the monotonic decrease of the continuous energy, we prove a crucial discrete counterpart: 
the fully discrete energy obtained from the convex splitting solution is non-increasing, i.e., it exhibits unconditional discrete energy stability.
\begin{theorem}\label{thm:energy-stability} 
    Assume that $0< u^0<1$. For any time step size $\tau>0$, the discrete energy of convex splitting scheme \eqref{scheme:cs} is dissipated with time steps, i.e.
    \begin{equation}
        E_h(u^{n+1})\leq E_h(u^n).
    \end{equation}
\end{theorem}
\begin{proof}
    The idea for this proof was inspired by the approach taken in \cite{eyre1998unconditionally}. 
    We decompose $F(u)$ into $F(u)=F_1(u)-F_2(u)$ where $F_1(u)=u\log(u)+(1-u)\log(1-u)$ and  $F_2(u)=\theta(u^2-u)$ are both convex. 
    From \eqref{eq:opt}, $J(u^{n+1})\leq J(u^n)$, and a simple fact that $\langle v,w \rangle_h = \langle vw,1 \rangle_h$ for $v,w \in \ch$, we have
    \begin{equation}
        \begin{aligned}
        &\frac{1}{2\tau}\|u^{n+1}-u^n\|_h^2+(\frac{\epsilon^2}{2}\|\nabla_hu^{n+1}\|_h^2+\langle F_1(u^{n+1}),1 \rangle_h)
        \\
        \leq& \langle u^{n+1}-u^n, F_2'(u^n)\rangle_h
        +(\frac{\epsilon^2}{2}\|\nabla_hu^{n}\|_h^2+\langle F_1(u^{n}),1 \rangle_h)
        ,
    \end{aligned}
    \end{equation}
    
the convexity of $F_2(u)$ implies\begin{equation}
    \langle u^{n+1}-u^n, F_2'(u^n)\rangle_h\leq 
    \langle F_2(u^{n+1})-F_2(u^n),1 \rangle_h.
\end{equation}
    After simplification, we arrive at
\begin{equation}
    E_h(u^{n+1})+\frac{1}{2\tau}\|u^{n+1}-u^n\|_h^2\leq E_h(u^n).
\end{equation}
\end{proof}


\subsection{Unconditional linear convergence}
% While the previous analysis guarantees the existence of a unique discrete solution of convex splitting scheme \eqref{scheme:cs}, 
% a practical question remains: can we compute it reliably? 
% This directs our focus to the iterative solver \Cref{alg:ADMM}. 
We begin by establishing the unconditional convergence of the algorithm, after which we prove that it converges at a linear rate.
\par The following lemma states the existence and uniqueness of the stationary point of the Lagrange function $\mathcal{L}$. 
\begin{lemma}
    Assume that $0<u^n<1$ and $u^n \in \ch$. For any time step size $\tau>0$, there exists a unique stationary point $(u_1^*,u_2^*,u_3^*)\in \ch \times \ch \times \ch$ of the Lagrange function $\mathcal{L}$ defined in \eqref{eq:Lagrange}.
\end{lemma}
\begin{proof}
    The stationary points are found by setting the partial derivatives of the function with respect to $u_1,u_2,u_3$ to zero, 
    which gives the following KKT system,
    \begin{subequations}\label{eq:stationary}
        \begin{align}
        &\frac{\partial \mathcal{L}}{\partial u_1}(u_1^*,u_2^*;u_3^*)
        =\frac{1}{\tau}(u_1^*-u^n)-\epsilon^2\Delta_hu_1^*+u_3^*=0,\\
        &\frac{\partial \mathcal{L}}{\partial u_2}(u_1^*,u_2^*;u_3^*)
        =\log(u_2^*)-\log(1-u_2^*)+\theta(1-2u^n)-u_3^*=0,\\
        &\frac{\partial \mathcal{L}}{\partial u_3}
        (u_1^*,u_2^*;u_3^*)
        =u_1^*-u_2^*=0.
        \end{align}
    \end{subequations}
    The above KKT system reduces to
    \begin{equation}
        \frac{u_1^*-u^n}{\tau}-\epsilon^2\Delta_hu_1^*
    +\log(u_1^*)-\log(1-u_1^*)+\theta(1-2u^n)=0,
    \end{equation}
    which is exactly obtained by substituting $u_1^*$ for $u^{n+1}$ in \eqref{scheme:cs}. By \Cref{thm:well-posedness}, we arrive at $u_1^*=u^{n+1}$, uniquely. 
    Define $(u_1^*,u_2^*,u_3^*)\in \ch \times \ch \times \ch$ as
    \begin{subequations}
        \begin{align}
            &u_1^*=u^{n+1},\\
            &u_2^*=u^{n+1},\\
            &u_3^*=-\frac{1}{\tau}(u_1^*-u^n)+\epsilon^2\Delta_hu_1^*=\log(u_2^*)-\log(1-u_2^*)+\theta(1-2u^n).
        \end{align}
    \end{subequations}
    Then $(u_1^*,u_2^*,u_3^*)$ satisfies the system \eqref{eq:stationary}, which means it is the unique stationary point of the Lagrange function $\mathcal{L}$.
\end{proof}
To prove the convergence of \Cref{alg:ADMM}, the sequence generated by the iterative solver is denoted by $(u_1^{(k)},u_2^{(k)},u_3^{(k)})$ and we adopt the notation $(u_1^*,u_2^*,u_3^*)$ established in the previous lemma. 
We define the error sequence
\begin{equation}
    (e_1^{(k)},e_2^{(k)},e_3^{(k)})=(u_1^{(k)}-u_1^*,u_2^{(k)}-u_2^*,u_3^{(k)}-u_3^*),\qquad k\in \mathbb{N},
\end{equation}
and the sequence $\{\Phi^{(k)}\}$ as
\begin{equation}\label{def:Phi}
    \Phi^{(k)}=\frac{1}{\alpha\rho}\|e_3^{(k)}\|_h^2+\rho\|e_1^{(k)}\|_h^2+
    \eta\rho\|e_1^{(k)}-e_2 ^{(k)}\|_h^2,\qquad k\in \mathbb{N}.
\end{equation}
where $\eta=\max(1-\alpha,1-\alpha^{-1})$. The techniques of convergence analysis are adapted from \cites{neal2011distributed,chen2017note,li2025overcoming}. 
We introduce the following useful inequality.
\begin{lemma}\label{lemma:Phi}
With the notations above, the following inequality holds,
\begin{equation}
\begin{aligned}
    \Phi^{(k)}-\Phi^{(k+1)}\geq 
    &\alpha c(\alpha)\rho\|e_1^{(k+1)}-e_1^{(k)}\|_h^2\\+
    &c(\alpha)\rho\|e_1^{(k+1)}-e_2^{(k+1)}\|_h^2,
\end{aligned}
\end{equation}
for all $k \in \mathbb{N}$, where $c(\alpha)=\min(1,1+\alpha^{-1}-\alpha)=\begin{cases}
    1,\alpha \in (0,1]\\
    1+\alpha^{-1}-\alpha,\alpha \in (1,+\infty)
\end{cases}$.
\end{lemma}
\begin{proof}
From \eqref{def:Phi}, we have
\begin{equation}
    \begin{aligned}
    &\Phi^{(k)}-\Phi^{(k+1)}\\
    =&\frac{1}{\alpha \rho}(\|e_3^{(k)}\|_h^2-\|e_3^{(k+1)}\|_h^2)+\rho (\|e_1^{(k)}\|_h^2-\|e_1^{(k+1)}\|_h^2)\\
    &+\eta \rho (\|e_1^{(k)}-
    e_2^{(k)}\|_h^2-\|e_1^{(k+1)}-
    e_2^{(k+1)}\|_h^2)\\
    =&(-2\langle e_3^{(k)},e_1^{(k+1)}-e_2^{(k+1)}\rangle_h-\alpha \rho \|e_1^{(k+1)}-e_2^{(k+1)}\|_h^2)\\
    &+\rho (\|e_1^{(k+1)}-e_1^{(k)}\|_h^2
    +2\langle e_1^{(k)}-e_1^{(k+1)},e_1^{(k+1)}\rangle_h
    )\\
    &+\eta \rho (\|e_1^{(k)}-
    e_2^{(k)}\|_h^2-\|e_1^{(k+1)}-
    e_2^{(k+1)}\|_h^2)
\end{aligned}
\end{equation}
From \eqref{eq:u1}, \eqref{eq:u2}, \eqref{eq:u3}, and the KKT system \eqref{eq:stationary}, we get the error equations,
\begin{subequations}
    \begin{align}
        &g(u_2^{(k+1)})-g(u_2^*)-\rho(e_1^{(k)}-e_2^{(k+1)})-e_3^{(k)}=0,
        \label{eq:e2}
        \\
        &(\frac{1}{\tau}-\epsilon^2\Delta_h)e_1^{(k+1)}
        +\rho(e_1^{(k+1)}-e_2^{(k+1)})+e_3^{(k)}=0,
        \label{eq:e1}
        \\
        &e_3^{(k+1)}=e_3^{(k)}+\alpha \rho(e_1^{(k+1)}-e_2^{(k+1)}),
        \label{eq:e3}
    \end{align}
\end{subequations}
where $g(u)=\log(u)-\log(1-u)$. \eqref{eq:e1} implies
\begin{equation}
    \langle (\frac{1}{\tau}-\epsilon^2\Delta_h)e_1^{(k+1)}
        +\rho(e_1^{(k+1)}-e_2^{(k+1)})+e_3^{(k)}, e_1^{(k+1)} \rangle_h=0,
\end{equation}
The positive definiteness of the operator $(\frac{1}{\tau}-\epsilon^2\Delta_h)$ leads to
\begin{equation}\label{eq:inequality1}
    \langle e_3^{(k)},e_1^{(k+1)}\rangle_h 
    \leq
    -\rho \langle e_1^{(k+1)}-e_2^{(k+1)},e_1^{(k+1)}\rangle_h.
\end{equation}
\eqref{eq:e2} implies
\begin{equation}
    \langle g(u_2^{(k+1)})-g(u_2^*)-\rho(e_1^{(k)}-e_2^{(k+1)})-e_3^{(k)},u_2^{(k+1)}-u_2^* \rangle_h=0
\end{equation}
The monotonicity property of $g(u)$ yields
\begin{equation}
    \langle g(u_2^{(k+1)})-g(u_2^*),e_2^{(k+1)} \rangle_h=
    \langle \rho(e_1^{(k)}-e_2^{(k+1)})+e_3^{(k)},e_2^{(k+1)} \rangle_h\geq 0,
\end{equation}
which reduces to 
\begin{equation}\label{eq:inequality2}
    \langle e_3^{(k)},e_2^{(k+1)} \rangle _h\geq -\rho 
    \langle e_1^{(k)}-e_2^{(k+1)},e_2^{(k+1)} \rangle_h
\end{equation}
It follows from \eqref{eq:inequality1} and \eqref{eq:inequality2} that
\begin{equation}
    -2\langle e_3^{(k)},e_1^{(k+1)}-e_2^{(k+1)}\rangle_h\geq 2\rho \langle e_1^{(k+1)}-e_2^{(k+1)},e_1^{(k+1)}\rangle_h - 2\rho
    \langle e_1^{(k)}-e_2^{(k+1)},e_2^{(k+1)} \rangle_h,
\end{equation}
which leads to
\begin{equation}
\begin{aligned}
&\Phi^{(k)}-\Phi^{(k+1)}\\
\geq &
2\rho \langle e_1^{(k+1)}-e_2^{(k+1)},e_1^{(k+1)}\rangle_h - 2\rho
    \langle e_1^{(k)}-e_2^{(k+1)},e_2^{(k+1)} \rangle_h \\
    &-\alpha \rho \|e_1^{(k+1)}-e_2^{(k+1)}\|_h^2\\
&+\rho (\|e_1^{(k+1)}-e_1^{(k)}\|_h^2
    +2\langle e_1^{(k)}-e_1^{(k+1)},e_1^{(k+1)}\rangle_h
    )\\
    &+\eta \rho (\|e_1^{(k)}-
    e_2^{(k)}\|_h^2-\|e_1^{(k+1)}-
    e_2^{(k+1)}\|_h^2)\\
=&(2-\alpha-\eta)\rho \|e_1^{(k+1)}-
    e_2^{(k+1)}\|_h^2 +
    2\rho \langle e_1^{(k+1)}-e_2^{(k+1)},e_1^{(k)}-e_1^{(k+1)} \rangle_h\\
&+\rho \|e_1^{(k+1)}-e_1^{(k)}\|_h^2+
\eta \rho \|e_1^{(k)}-
    e_2^{(k)}\|_h^2
\end{aligned}
\end{equation}
Taking $k$-th and $k+1$-th in \eqref{eq:e1}, we obtain that
\begin{equation}
    (\tau^{-1}-\epsilon^2 \Delta_h)(e_1^{(k+1)}-e_1^{(k)})+\rho (e_1^{(k+1)}-e_2^{(k+1)}-(e_1^{(k)}-e_2^{(k)}))+(e_3^{(k)}-e_3^{(k-1)})=0.
\end{equation}
We take the inner product with $e_1^{(k+1)}-e_1^{(k)}$. Then, by exploiting the positive-definiteness of $\tau^{-1}-\epsilon^2 \Delta_h$ 
and the definition of $e_3^{(k)}$, we infer the following useful inequality that
\begin{equation}\label{eq:inequality4}
    \langle e_1^{(k+1)}-e_2^{(k+1)},e_1^{(k+1)}-e_1^{(k)}\rangle_h
    \leq(1-\alpha) \langle e_1^{(k)}-e_2^{(k)},
    e_1^{(k+1)}-e_1^{(k)}\rangle_h,
\end{equation}
which implies
\begin{equation}
\begin{aligned}
    &\Phi^{(k)}-\Phi^{(k+1)}\\
\geq & (2-\alpha-\eta)\rho \|e_1^{(k+1)}-
    e_2^{(k+1)}\|_h^2-2(1-\alpha)\rho \langle e_1^{(k)}-e_2^{(k)},
    e_1^{(k+1)}-e_1^{(k)}\rangle_h\\
    &+\rho \|e_1^{(k+1)}-e_1^{(k)}\|_h^2
    +\eta \rho \|e_1^{(k)}-
    e_2^{(k)}\|_h^2
\end{aligned}
\end{equation}

If $\alpha \in (0,1]$, $\eta=1-\alpha$. From the inequality $2ab\leq a^2+b^2$, we derive that
\begin{equation}\label{eq:a<1}
    \begin{aligned}
    &\Phi^{(k)}-\Phi^{(k+1)}\\
\geq & (2-\alpha-(1-\alpha))\rho \|e_1^{(k+1)}-
    e_2^{(k+1)}\|_h^2\\
    &-(1-\alpha)\rho (\| e_1^{(k)}-e_2^{(k)}\|_h^2+
    \|e_1^{(k+1)}-e_1^{(k)}\|_h^2)\\
    &+\rho \|e_1^{(k+1)}-e_1^{(k)}\|_h^2
    +(1-\alpha) \rho \|e_1^{(k)}-
    e_2^{(k)}\|_h^2\\
    =&\alpha \rho \|e_1^{(k+1)}-e_1^{(k)}\|_h^2+\rho \|e_1^{(k+1)}-
    e_2^{(k+1)}\|_h^2.
\end{aligned}
\end{equation}
If $\alpha \in (1,+\infty)$, $\eta=1-\alpha^{-1}$. It follows from the inequality $2ab\leq \alpha^{-1}a^2+\alpha b^2$ that
\begin{equation}\label{eq:a>1}
\begin{aligned}
&\Phi^{(k)}-\Phi^{(k+1)}\\
\geq &(2-\alpha-(1-\alpha^{-1}))\rho \|e_1^{(k+1)}-
    e_2^{(k+1)}\|_h^2\\
    &+(\alpha-1)\rho (-\alpha^{-1}\| e_1^{(k)}-e_2^{(k)}\|_h^2-
    \alpha \|e_1^{(k+1)}-e_1^{(k)}\|_h^2)\\
    &+\rho \|e_1^{(k+1)}-e_1^{(k)}\|_h^2
    +(1-\alpha^{-1}) \rho \|e_1^{(k)}-
    e_2^{(k)}\|_h^2\\
    =&(1+\alpha-\alpha^2)\rho \|e_1^{(k+1)}-e_1^{(k)}\|_h^2+(1-\alpha+\alpha^{-1})\rho \|e_1^{(k+1)}-
    e_2^{(k+1)}\|_h^2.
\end{aligned}
\end{equation}
The proof is completed by the above two inequalities \eqref{eq:a<1} and \eqref{eq:a>1}.
\end{proof}
Armed with the lemma above \ref{lemma:Phi}, we proceed to prove unconditional convergence.
\begin{theorem}\label{thm:convergence}
    Assume that the step size $\alpha \in (0,\frac{\sqrt{5}+1}{2})$ for multiplier update. Then the iterative solver \Cref{alg:ADMM} converges,
    \begin{equation}
        \lim_{k\to +\infty} u_1^{(k)}=
        \lim_{k\to +\infty} u_2^{(k)}=u^{n+1}.
    \end{equation}
\end{theorem}
\begin{proof}
    Notice that $c(\alpha)>0$ is equivalent to $\alpha \in (0,\frac{\sqrt{5}+1}{2})$. By lemma \ref{lemma:Phi}, the sequence $\Phi^{(k)}$ is bounded and convergent because it is monotonically decreasing and bounded below by $0$ \cite{1976Principles}, thus the sequences
    \begin{align*}
        \{e_3^{(k)}\},\{e_1^{(k)}\},\{e_1^{(k)}-e_2^{(k)}\}
    \end{align*}
    are bounded. Then, the sequences
       $ \{u_1^{(k)}\},\{u_2^{(k)}\},\{u_3^{(k)}\}$ are bounded. Therefore, there exists a subsequence $\{(u_1^{(k_l)},u_2^{(k_l)},u_3^{(k_l)}) \}$ which converges to the unique stationary point of $\mathcal{L}$, i.e.
       \begin{equation}
           \lim_{l\to +\infty}(u_1^{(k_l)},u_2^{(k_l)},u_3^{(k_l)})=(u_1^*,u_2^*,u_3^*)=(u^{n+1},u^{n+1},-\frac{1}{\tau}(u_1^*-u^n)+\epsilon^2\Delta_hu_1^*).
       \end{equation}
      Since $\{\Phi^{(k)}\}$ is convergent and has a subsequence $\{\Phi^{(k_l)}\}$ converging to $0$, the sequence $\{\Phi^{(k)}\}$ converges to $0$ itself \cite{1976Principles}. Then we have
      \begin{equation}
          \lim_{k\to +\infty}\|e_3^{(k)}\|_h=\lim_{k\to +\infty}\|e_1^{(k)}\|_h=\lim_{k\to +\infty}\|e_1^{(k)}-e_2^{(k)}\|_h=0.
      \end{equation}
      The triangle inequality \begin{equation}
          \|e_2^{(k)}\|_h\leq \|e_2^{(k)}-e_1^{(k)}\|_h+\|e_1^{(k)}\|_h 
      \end{equation}
      implies 
      \begin{equation}
          \lim_{k\to +\infty}\|e_2^{(k)}\|_h=0. 
      \end{equation}
      We have thus completed the proof.
\end{proof}

In what follows, we prove that the algorithm converges at a linear rate. 
The proof relies on the following lemma, which provides a sufficient condition for linear convergence. 
\begin{lemma}\label{lm:J}
    The objective function $J(u)=J_1(u)+J_2(u)$ in \eqref{eq:Ju} satisfies the regularity conditions required for linear convergence: 
$J_1(u)$ is strongly convex with a Lipschitz continuous gradient, and $J_2(u)$ is strongly convex.
\end{lemma}
\begin{proof}
    By direct computation, one obtains
\begin{subequations}
\begin{align}
\label{eq:Hessian1}
\nabla^2J_1(u)=\tau^{-1}I-\epsilon^2\Delta_h,\\
\label{eq:Hessian2}
\nabla^2J_2(u)=G=\text{diag}\{\frac{1}{u_j(1-u_j)}\}_{j=0}^{N^2-1},
\end{align}
\end{subequations}
From \eqref{eq:Hessian1}, the function $J_1(u)$ is strongly convex with constant $\mu_1>0$ 
and has a Lipschitz continuous gradient with constant $L_1>0$, which are given by
\begin{align}
    \mu_1=\tau^{-1},\\
    L_1=\tau^{-1}+\frac{8\epsilon^2}{h^2}.
\end{align}
This is due to the known spectral properties of the negative discrete Laplacian matrix $-\Delta_h$ 
with $\lambda_{\min}(-\Delta_h)=0$ and $\lambda_{\max}(-\Delta_h)=\frac{8}{h^2}$.
\par Similarly, with \eqref{eq:Hessian2}, the function $J_2(u)$ is strongly convex with constant $\mu_2=4>0$, 
yet its gradient is not Lipschitz continuous.
\end{proof}
Defining the sequence 
\begin{equation}
R^{(k)}=\rho \|e_1^{k}\|_h^2+\rho^{-1}\|e_3^{k}\|_h^2, 
\end{equation}
the following theorem describes the linear convergence rate of $\{R^{(k)}\}$ in \Cref{alg:ADMM}.
\begin{theorem}
Let the sequence $\{R^{(k)}\}$ be defined as above. It is linearly convergent in the sense that there extists a constant $\delta>0$ such that 
\begin{equation}
    R^{(k)}\geq (1+\delta) R^{(k+1)},
\end{equation}
for all $k\in \mathbb{N}$. In the special case where the step size for multiplier update is fixed as $\alpha=1$, $\delta$ takes the following explicit form:
\begin{equation}
\delta=2(\frac{\rho}{\mu_1}+\frac{L_1}{\rho})^{-1}.
\end{equation}
Choosing $\rho=\sqrt{L_1\mu_1}$ yields the largest $\delta$,
\begin{equation}
\delta_{\max}=\sqrt{\frac{\mu_1}{L_1}}=\sqrt{\frac{h^2}{h^2+8\tau \epsilon^2}}.
\end{equation}
\end{theorem}
\begin{proof}
    \ref{lm:J} implies that our optimization problems falls into the scenario $1$ of problems analyzed in \cite{deng2016global}. 
According to \cite[Theorem 3.4]{deng2016global}, our iterative solver \Cref{alg:ADMM} consequently achieves a linear convergence rate. 
The explicit formula of convergence rate is given by \cite[Corollary 3.6]{deng2016global}.
\end{proof}

\begin{remark}
We briefly discuss the bound-preserving property of the algorithm. 
For the update of $u_2$, the presence of the logarithmic term ensures $u_2^{(k)}\in(0,1)$ at a point-wise level for all $k \in \mathbb{N}$. 
\Cref{thm:convergence} guarantees that $u_2^{(k)}$ is a good approximation of $u^{n+1}$ after a sufficiently large number of iterations. 
Consequently, we obtain the bound-preserving property.
\end{remark}





